\chapter{Introduction}\label{chap:introduction}

This chapter sets out the question the book answers and why it matters. It also explains how the rest of the book is organized, so a reader can choose where to start.

\section{Motivation}\label{sec:motivation}

Measuring uncertainty well is older than most of the methods used to do it. The foundations of information theory \citep{shannon1948} still shape how we reason about noise and signal, and resampling \citep{efron1979} made it practical to attach error bars to estimates that have no tidy formula.\footnote{Footnotes suit short asides that would interrupt the main argument.}

\section{Outline}\label{sec:outline}

\cref{chap:methods} describes the approach and the estimators it compares. \cref{chap:results} reports what happens when those estimators are run on simulated data. The appendix collects a derivation that the main text only sketches.
